\documentclass[10pt]{article}
\usepackage[letterpaper,margin=0.68in]{geometry}
\usepackage{amsmath,amssymb}
\usepackage{enumitem}
\usepackage{xcolor}
\usepackage{fancyhdr}
\usepackage{microtype}
\usepackage[most]{tcolorbox}
\definecolor{olive}{HTML}{4D5430}
\definecolor{gold}{HTML}{9A7B22}
\definecolor{pale}{HTML}{F5F1DC}
\definecolor{softgreen}{HTML}{EDF0E2}
\pagestyle{fancy}\fancyhf{}\lhead{\textcolor{olive}{MATH\& 107 Answer Key}}\rhead{\textcolor{gold}{Fall 2026}}\cfoot{\thepage}
\setlength{\headheight}{14pt}
\setlength{\parindent}{0pt}\setlength{\parskip}{5pt}
\setlist[enumerate]{leftmargin=*,itemsep=5pt,topsep=4pt}
\newtcolorbox{solbox}[1]{colback=pale,colframe=olive,boxrule=.7pt,arc=2mm,title=\textbf{#1},fonttitle=\color{olive},left=2mm,right=2mm,top=1.5mm,bottom=1.5mm}
\newcommand{\modulekey}[2]{\clearpage{\Large\bfseries\textcolor{olive}{Module #1: #2 — Answer Key}}\par\textcolor{gold}{\rule{\linewidth}{1.2pt}}\par}
\begin{document}
\modulekey{9}{Linear Programming}
\textbf{Prequiz}
\begin{enumerate}
\item Labor:
\[
\boxed{3x+2y\le120}.
\]
Land:
\[
\boxed{x+y\le50}.
\]
Nonnegativity: \(\boxed{x\ge0,y\ge0}\). Objective:
\[
\boxed{\text{maximize }P=200x+150y}.
\]
\item
\[
f(0,0)=0,\quad f(8,0)=40,\quad f(4,6)=68,\quad f(0,8)=64.
\]
Maximum \(=\boxed{68}\) at \(\boxed{(4,6)}\).
\item
\[
C(8,0)=24,\quad C(2,6)=36,\quad C(0,10)=50.
\]
Minimum \(=\boxed{24}\) at \(\boxed{(8,0)}\).
\item
\[
f(0,0)=0,\quad f(4,0)=12,\quad f(3,2)=15,\quad f(0,5)=15.
\]
Optimal vertices: \(\boxed{(3,2)}\) and \(\boxed{(0,5)}\). Every point on the edge joining them is also optimal.
\item Check nearby feasible \textbf{integer} points, because the continuous optimum uses fractional servings and is not directly allowable when servings must be whole.
\end{enumerate}

\textbf{Matching quiz}
\begin{enumerate}
\item Bread:
\[
\boxed{2x+y\le80}.
\]
Filling:
\[
\boxed{3x+4y\le150}.
\]
Also \(x,y\ge0\). Objective:
\[
\boxed{\text{maximize }R=4x+5y}.
\]
\item
\[
R(0,0)=0,\quad R(6,0)=36,\quad R(4,5)=59,\quad R(0,7)=49.
\]
Maximum \(=\boxed{59}\) at \(\boxed{(4,5)}\).
\item For \(C=5x+2y\):
\[
C(6,0)=30,\quad C(3,6)=27,\quad C(0,12)=24.
\]
Minimum \(=\boxed{24}\) at \(\boxed{(0,12)}\).
\item Every point on the edge joining the two tied adjacent vertices is optimal.
\item \(7.5\) machines is not physically possible if machines must be whole. Check nearby feasible integer points, such as points with \(x=7\) or \(x=8\), before making the final practical choice.
\end{enumerate}
\end{document}
